% sec-06: the advisory boundary, stop authority, verification.
% Figures 9, 10, 12, 13, 14, 15.
\section{Physical AI Governance}

\subsection{Three Boundaries}

\textbf{The model has no path to the arms.} Advisory output is rendered to a
display. No generated token reaches a motion controller, and no software in the
advisory path holds write access to the robot's command channel. The operating
surgeon approves every motion \cite{onpremwhippl}.

\textbf{Stopping is bounded and measured.} Arm-level stop is specified at 3
milliseconds and system-wide stop at 500 milliseconds. Both are measured on the
bench before the first participant is screened and re-measured at every cohort
boundary. A specification that is asserted rather than measured is not a
boundary; it is a hope.

\textbf{Every artifact is verified before use.} The model runs on premises,
pinned to a repository commit, and writes a timestamped record for every
recommendation. Conversion to open operation ordered for patient safety is
recorded as clinically appropriate; conversion caused by device, advisory, or
integration failure is counted separately against the feasibility endpoint
\cite{phase1ind}.

\begin{appfloat}[!tb]
\begin{appfig}[diagrams (python)-type / clustered topology]
% Three stacked clusters inside the boundary so the boundary rule can run full
% height without crossing one.  Tiles 2.4 apart, labels 5.4mm below tile
% centres.  The struck link marks the path from the model to the controller.
\dgnodew{dgtiled}{0}{0}{\glyphrobot}{Eight-arm robotic platform}{t1}
\dgnode{dgtile}{2.4}{0}{\glyphcpu}{Motion controller}{t2}
\dgnode{dgtile}{4.8}{0}{\glyphmon}{Advisory display}{t3}
\begin{scope}[on background layer]
\node[dgcluster,fit=(t1)(t1l)(t2)(t2l)(t3)(t3l),inner sep=7pt] (co) {};
\end{scope}
\node[dgctitle,anchor=south west] at ([yshift=1mm]co.north west) {Operative};
\dgnode{dgtile}{0}{-2.7}{\glyphai}{On-premises model, pinned commit}{t4}
\dgnode{dgtile}{2.4}{-2.7}{\glyphdb}{Local audit log}{t5}
\dgnode{dgtile}{4.8}{-2.7}{\glyphserver}{Offline model registry}{t6}
\begin{scope}[on background layer]
\node[dgcluster,fit=(t4)(t4l)(t5)(t5l)(t6)(t6l),inner sep=7pt] (ca) {};
\end{scope}
\node[dgctitle,anchor=south west] at ([yshift=1mm]ca.north west) {Advisory};
\dgnode{dgtile}{0}{-5.4}{\glyphlock}{Identity broker and access control}{t7}
\dgnode{dgtile}{2.4}{-5.4}{\glyphsignal}{Site monitoring endpoint}{t8}
\begin{scope}[on background layer]
\node[dgcluster2,fit=(t7)(t7l)(t8)(t8l),inner sep=7pt] (cc) {};
\end{scope}
\node[dgctitle2,anchor=south west] at ([yshift=1mm]cc.north west) {Control};
\dgnodeg{dgtileg}{9.4}{-2.7}{\glyphcloud}{Public archive, DOI minted}{o1}
\begin{scope}[on background layer]
\node[dgcluster2,fit=(o1)(o1l),inner sep=7pt] (cx) {};
\end{scope}
\node[dgctitle2,anchor=south west] at ([yshift=1mm]cx.north west) {Outside};
\draw[protoblack,line width=1pt] (7.2,1.2) -- (7.2,-6.3);
\node[font=\tiny\sffamily\bfseries,text=protoblack,rotate=90,anchor=south] at (7.05,-2.5) {trust boundary};
\draw[dgedgeb] (t4) -- (t3);
\draw[dgedgeb] (t4) -- (t5);
\draw[dgedgeb] (t3) -- node[font=\tiny\sffamily,pos=0.5,fill=protowhite,inner sep=1pt] {human} (t2);
\draw[dgedgeb] (t2) -- (t1);
\draw[dgedge] (t6) -- (t4);
\draw[dgedge] (t7) -- (t4);
\draw[dgedge] (t5) -- (t8);
\draw[dgedged] (t5.east) to[out=20,in=180,looseness=0.8] (o1.west);
\pxmark{2.4}{-1.35}
\node[font=\tiny\sffamily,text=protogray,anchor=west,text width=26mm,align=left]
  at (2.75,-1.35) {no path from the model to the controller};
\node[font=\tiny\sffamily,text=protogray,anchor=north,text width=44mm,align=center]
  at (9.4,-4.6) {Records leave. Nothing returns: there is no inbound edge across
  the boundary.};
\end{appfig}
\vspace{-0.7cm}
\figcaption{Nine components, one trust boundary, and one direction of travel.
Records leave for the public archive; nothing returns, and no path from the
model reaches a motion controller at any point.}
\end{appfloat}

\subsection{The Operative Day}

Two figures are needed for one operative day, and neither implies the other. The
sequence gives the conversation, message by message, which is what an operating
team recognises. The state machine gives the guards, which is what a reviewer
checks. A sequence with no guards cannot say what is forbidden; a state machine
with no messages cannot say who speaks.

\begin{appfloat}[!tb]
\begin{appfig}[mermaid-type / sequence diagram]
% Five lifelines 2.9 apart; eleven messages 0.62 apart.  Every label sits on a
% white ground so a message never overprints a lifeline.
\foreach \i/\x/\lab in {1/0/{Operating surgeon}, 2/2.9/{Console and controller},
  3/5.8/{Eight-arm platform}, 4/8.7/{On-premises model}, 5/11.6/{Audit log}}{
  \node[mmactor,text width=20mm] (a\i) at (\x,0) {\lab};
  \draw[mmlife] (\x,-0.45) -- (\x,-7.35);}
\foreach \y/\f/\t/\msg in {%
  -0.90/0/2.9/{1. set operative plan for the step},
  -1.52/2.9/8.7/{2. telemetry snapshot, read only},
  -2.14/8.7/0/{3. advisory text rendered to a display},
  -2.76/8.7/11.6/{4. recommendation written with timestamp},
  -3.38/0/2.9/{5. approve, modify, or reject},
  -4.00/2.9/5.8/{6. motion command},
  -4.62/5.8/2.9/{7. force and position feedback},
  -5.24/0/11.6/{8. decision written with latency},
  -6.10/0/2.9/{9. request stop},
  -6.72/2.9/5.8/{10. arm stop 3 ms, system stop 500 ms}}{
  \draw[mmmsg] (\f,\y) -- (\t,\y)
    node[midway,above,font=\tiny\sffamily,fill=protowhite,inner sep=1.2pt] {\msg};}
\draw[pagrayd,dashed,line width=0.4pt] (-0.9,-5.65) -- (12.5,-5.65);
\pxmark{7.25}{-5.65}
\node[font=\tiny\sffamily,text=protogray,anchor=west,text width=44mm,align=left]
  at (7.6,-5.65) {no message exists from the model to the console or to the platform};
\end{appfig}
\vspace{-0.7cm}
\figcaption{One operative day as a message sequence. Every arrow reaching the
robot leaves a human first, and the two audit writes are the only messages that
cross the hospital trust boundary.}
\end{appfloat}

\begin{appfloat}[!tb]
\begin{appfig}[plantuml-type / state diagram with guards]
% Four normal states on the upper rank at y = 0, pitch 3.3.  Two terminal-side
% states on the lower rank at y = -2.4.  The two human-only transitions leave
% one state, so they are drawn at bend left=14 and bend right=14 with equal
% looseness; any larger bend re-enters the Composing box.
\node[umlinit] (i) at (-1.5,0) {};
\node[umlstatesoft,text width=24mm] (s1) at (0.3,0) {Idle};
\node[umlstatesoft,text width=24mm] (s2) at (3.6,0) {Reading telemetry};
\node[umlstatesoft,text width=24mm] (s3) at (6.9,0) {Composing};
\node[umlstate,text width=24mm] (s4) at (10.2,0) {Rendered};
\node[umlstateon,text width=24mm] (s5) at (10.2,-2.4) {Accepted};
\node[umlstategray,text width=24mm] (s6) at (6.9,-2.4) {Rejected};
\node[umlstatesoft,text width=24mm] (s7) at (3.6,-2.4) {Logged};
\node[umlstategray,text width=24mm] (s8) at (0.3,-2.4) {Degraded};
\node[umlfinal] (f) at (-1.5,-2.4) {};
\draw[umlarrow] (i) -- (s1);
\draw[umlarrow] (s1) -- node[umlguard,pos=0.5] {[case open]} (s2);
\draw[umlarrow] (s2) -- node[umlguard,pos=0.5] {[snapshot complete]} (s3);
\draw[umlarrow] (s3) -- node[umlguard,pos=0.5] {[template check passes]} (s4);
\draw[umlarrow] (s4) to[bend left=14] node[umlguard,pos=0.5] {[surgeon accepts]} (s5);
\draw[umlarrow] (s4.south west) to[bend right=14] node[umlguard,pos=0.45] {[surgeon rejects]} (s6.north east);
\draw[umlarrow] (s5) -- (s7);
\draw[umlarrow] (s6) -- (s7);
\draw[umlarrow] (s7) -- node[umlguard,pos=0.5] {[step closed]} (s8);
\draw[umlarrow] (s8) -- (f);
\draw[umldash] (s2.south) -- node[umlguard,pos=0.55] {[telemetry stale $>$ 250 ms]} (s8.north east);
\draw[umldash] (s3.south) -- node[umlguard,pos=0.5] {[model latency $>$ 2 s]} (s8.north);
\node[umlbar,minimum width=3mm,minimum height=1.4mm] at (10.9,-0.95) {};
\node[umlbar,minimum width=3mm,minimum height=1.4mm] at (8.9,-1.25) {};
\node[font=\tiny\itshape,text=protogray,anchor=north,text width=110mm,align=center]
  at (5.2,-3.3) {The two marked transitions are the only ones a human can fire,
  and they are the two that separate a rendered recommendation from an executed
  motion. Degraded renders nothing: it is silence, not a fallback mode.};
\end{appfig}
\vspace{-0.7cm}
\figcaption{Six advisory states and the guard on every transition. Two guards
can only be satisfied by a human action, and they are the two that separate a
rendered recommendation from an executed motion.}
\end{appfloat}

\subsection{How a Failure Propagates}

\begin{appfloat}[!tb]
\begin{appfig}[graphviz-type / fault tree]
% One top event, one OR gate, four branches, four gates, eight leaves.  Ranks at
% y = 0, -1.05, -2.4, -3.45, -4.9.  Every edge is vertical or near-vertical and
% none crosses a gate glyph.
\node[gvboxk,text width=44mm] (top) at (5.6,0) {Unintended motion reaches the patient};
\umlgateor{5.6}{-1.05}
\draw[gvedge] (5.6,-0.79) -- (top.south);
\foreach \i/\x/\lab in {1/0.4/{Command path fails}, 2/3.9/{Stop path fails},
  3/7.4/{Advisory mistaken for a command}, 4/10.9/{Wrong step authorised}}{
  \node[gvboxs,text width=25mm] (b\i) at (\x,-2.4) {\lab};}
\draw[gvedge] (b1.north) -- (5.28,-1.35);
\draw[gvedge] (b2.north) -- (5.44,-1.38);
\draw[gvedge] (b3.north) -- (5.76,-1.38);
\draw[gvedge] (b4.north) -- (5.92,-1.35);
\umlgateand{0.4}{-3.45}
\umlgateand{3.9}{-3.45}
\umlgateand{7.4}{-3.45}
\umlgateor{10.9}{-3.45}
\foreach \i/\x in {1/0.4, 2/3.9, 3/7.4, 4/10.9}{\draw[gvedge] (\x,-3.19) -- (b\i.south);}
\foreach \i/\x/\a/\b in {1/0.4/{Controller fault}/{Interlock not armed},
  2/3.9/{Arm stop $>$ 3 ms}/{System stop $>$ 500 ms},
  3/7.4/{Display omits the advisory banner}/{Surgeon does not confirm},
  4/10.9/{Plan and telemetry disagree}/{No second confirmation}}{
  \node[gvboxg,text width=21mm] (l\i a) at (\x-1.2,-4.9) {\a};
  \node[gvboxg,text width=21mm] (l\i b) at (\x+1.2,-4.9) {\b};
  \draw[gvedge] (l\i a.north) -- (\x-0.28,-3.72);
  \draw[gvedge] (l\i b.north) -- (\x+0.28,-3.72);}
\node[font=\tiny\itshape,text=protogray,anchor=north,text width=118mm,align=center]
  at (5.6,-5.7) {Three branches are AND gates, so no single failure reaches the
  patient. The fourth is an OR gate and both of its leaves are procedural, which
  is the argument: hardware redundancy does not cover a procedural gap.};
\end{appfig}
\vspace{-0.7cm}
\figcaption{One top event and every path to it. Three of the four branches are
AND gates, so no single failure reaches the patient, and the one OR branch is
the one governed by procedure rather than by hardware.}
\end{appfloat}

\subsection{What Is Released}

\begin{appfloat}[!tb]
\begin{appfig}[d2-type / sql tables]
% Participant centred above; five dependents fanned below at two levels, so no
% join edge passes under a table.  operative_step sits directly beneath
% participant because it carries the schema's only two-hop path.
\node[d2sql,text width=32mm] (pt) at (5.6,1.6)
  {\textbf{participant}\nodepart{two}participant\_id : uuid PK\\dose\_level : int\\kras\_variant : text\\resectability : text};
\node[d2sql,text width=32mm] (os) at (5.6,-1.5)
  {\textbf{operative\_step}\nodepart{two}step\_id : uuid PK\\participant\_id : uuid FK\\step\_ordinal : int\\arm\_ids : int array};
\node[d2sql,text width=30mm] (ad) at (0.4,-1.5)
  {\textbf{advisory\_decision}\nodepart{two}decision\_id : uuid PK\\step\_id : uuid FK\\accepted : bool\\latency\_ms : float};
\node[d2sql,text width=30mm] (se) at (10.8,-1.5)
  {\textbf{stop\_event}\nodepart{two}stop\_id : uuid PK\\step\_id : uuid FK\\measured\_ms : float\\scope : text};
\node[d2sql,text width=30mm] (ae) at (0.4,-4.6)
  {\textbf{adverse\_event}\nodepart{two}ae\_id : uuid PK\\participant\_id : uuid FK\\attribution : enum\\ctcae\_grade : int};
\node[d2sql,text width=30mm] (sp) at (10.8,-4.6)
  {\textbf{specimen}\nodepart{two}specimen\_id : uuid PK\\participant\_id : uuid FK\\pathway\_pd : float\\response\_grade : text};
\draw[d2edgeb] (os.north) -- (pt.south);
\draw[d2edge] (ad.east) -- (os.west);
\draw[d2edge] (se.west) -- (os.east);
\draw[d2edge] (ae.north) to[out=90,in=180,looseness=0.7] (pt.west);
\draw[d2edge] (sp.north) to[out=90,in=0,looseness=0.7] (pt.east);
\node[font=\tiny\itshape,text=protogray,anchor=north,text width=112mm,align=center]
  at (5.6,-6.4) {Every released record traces to one participant and one
  operative step. That is what makes an independent re-analysis possible without
  access to the site.};
\end{appfig}
\vspace{-0.7cm}
\figcaption{Six tables and the five keys that join them. Every released record
can be traced to one participant and one operative step, which is what makes an
independent re-analysis of the cohort possible at all.}
\end{appfloat}

\subsection{Who Produced Which Artifact}

The documents this paper summarises were themselves produced by a small set of
frontier models with separated roles, and separating them is what makes the
review chain meaningful: no model reviews its own output
\cite{repo2026}.

\begin{apptable}
\begin{tabularx}{\textwidth}{>{\raggedright\arraybackslash}p{2.9cm} >{\raggedright\arraybackslash}p{2.4cm} >{\raggedright\arraybackslash}X}
\headrow \thead{Model} & \thead{Role} & \thead{Artifacts produced}\\
\midrule
Claude Code & Author & Investigational new drug applications; trial protocols; FDA, ICH and NIH regulation adaptions; legislation revisions; machine-readable diagrams; repository platforms \\
ChatGPT Codex & First reviewer & Verification, validation and uncertainty quantification of author outputs; meta-analyses of trial papers; external standards; PDF reading and writing \\
Google Gemini & Second reviewer & Second peer review; accelerated code problem solving; meta-verification of prior stages; administrative automation \\
\end{tabularx}
\end{apptable}
\tabcap{The three frontier-model roles, taken verbatim from the repository's own
record. The separation is the point: the author model does not review itself,
and the second reviewer checks the first.}

\begin{appfloat}[!tb]
\begin{appfig}[diagrams (python)-type / clustered by vendor]
% Three vendor clusters on three ranks 3.0 apart, three tiles each at x = 0,
% 2.6, 5.2.  Two return edges at bend left=22 keep clear of the tile field.
\dgnode{dgtile}{0}{0}{\glyphdoc}{IND and protocols}{v1}
\dgnode{dgtile}{2.6}{0}{\glyphshield}{Regulation adaptions}{v2}
\dgnode{dgtile}{5.2}{0}{\glyphchart}{Machine-readable diagrams}{v3}
\begin{scope}[on background layer]
\node[dgcluster,fit=(v1)(v1l)(v2)(v2l)(v3)(v3l),inner sep=7pt] (k1) {};
\end{scope}
\node[dgctitle,anchor=south west] at ([yshift=1mm]k1.north west) {A. Claude Code, author};
\dgnode{dgtilem}{0}{-3.0}{\glyphgear}{VVUQ of author outputs}{w1}
\dgnode{dgtilem}{2.6}{-3.0}{\glyphflask}{Trial meta-analyses}{w2}
\dgnode{dgtilem}{5.2}{-3.0}{\glyphdoc}{External standards, PDF}{w3}
\begin{scope}[on background layer]
\node[dgcluster,fit=(w1)(w1l)(w2)(w2l)(w3)(w3l),inner sep=7pt] (k2) {};
\end{scope}
\node[dgctitle,anchor=south west] at ([yshift=1mm]k2.north west) {B. ChatGPT Codex, first reviewer};
\dgnodeg{dgtileg}{0}{-6.0}{\glyphteam}{Second peer review}{x1}
\dgnodeg{dgtileg}{2.6}{-6.0}{\glyphcpu}{Code problem solving}{x2}
\dgnodeg{dgtileg}{5.2}{-6.0}{\glyphlink}{Meta-verification}{x3}
\begin{scope}[on background layer]
\node[dgcluster2,fit=(x1)(x1l)(x2)(x2l)(x3)(x3l),inner sep=7pt] (k3) {};
\end{scope}
\node[dgctitle2,anchor=south west] at ([yshift=1mm]k3.north west) {C. Google Gemini, second reviewer};
\draw[dgedgeb] (k1.south) -- (k2.north);
\draw[dgedgeb] (k2.south) -- (k3.north);
\draw[dgedged] (k3.east) to[bend left=22] node[font=\tiny\sffamily,pos=0.5,fill=protowhite,inner sep=1pt] {findings return} (k1.east);
\draw[dgedged] (k3.east) to[bend left=16] (k2.east);
\node[font=\tiny\itshape,text=protogray,anchor=north,text width=100mm,align=center]
  at (2.6,-7.3) {The chain is a cycle rather than a line, because a review
  finding that is recorded but not acted on is not a review.};
\end{appfig}
\vspace{-0.7cm}
\figcaption{Three vendors, nine roles, and the single review chain that connects
them. No model reviews its own output, and the chain is drawn as a cycle because
the second reviewer's finding returns to the author.}
\end{appfloat}
